
Mới có mấy tháng sau sự ra đi của PGS Văn Như Cương, mà như lúc sinh thời ông thường ví von, đó là sự “ánh xạ” của một phần tử  từ không gian này đến một không gian khác không “đẳng cấu”, giờ lại đến lượt GS Phan Đình Diệu ! Buồn biết bao nhiêu !

Cũng như Văn Như Cương và nhiều trí thức chân chính khác, hệ tư duy trong con người Phan Đình Diệu là hoàn toàn minh bạch. Từ đó, ông nhìn mọi hệ thống lành mạnh phải là hệ thống minh bạch, trong đó mỗi phần tử đều có chức năng rõ ràng và hệ thống được vận hành bằng những mối quan hệ cũng rõ ràng ! Mọi phần tử của hệ thống đó đều hiểu và hành động nhất quán và trung thực theo các chức năng và quan hệ đó. 
Không chấp nhận sự mập mờ khái niệm và sự “vận dụng” tuỳ tiện.
Ở đâu cũng thế và lúc nào cũng thế ! 

Ngành Toán học mà Phan Đình Diệu đã nghiên cứu và phát triển là Toán học kiến thiết. Nhà toán học tìm ra được điều gì mới thì phải chứng minh nó bằng cách chỉ ra được một quy trình (một thuật toán) để đi đến điều phát minh đó ! 
Phương pháp phản chứng ( tức là chỉ ra rằng nếu công nhận ngược lại thì sẽ đi đến mâu thuẫn), mặc nhiên thừa nhận luật bài trung ( loại trừ cái thứ ba ), chứ không hề dẫn dắt người ta đến điều cần chứng minh. 

Vì thế, mà Toán học kiến thiết không chấp nhận Phương pháp phản chứng ! Cũng tức là không công nhận Luật bài trung !
Tư tưởng đó mới mẻ biết chừng nào ! Và cần thiết biết chừng nào cho việc lập trình mà sự nở rộ và phát triển như vũ bão của nó đã dẫn đến quá trình tự động hoá, không phải chỉ trong vận hành cơ điện vật lý…mà cả trong và chủ yếu là trong hệ thống tự động hoá điều khiển, trong tư duy sáng tạo và trí tuệ nhân tạo !

Bây giờ, khi mà cuộc cách mạng 4.0 đang đến gần thì tư duy đó được công nhận là  sáng tạo và thực tiễn, nhưng 40, 50 năm trước, nó xa vời biết chừng nào, nhất là đối với những người mà tư duy đã dừng lại với “Đại công trường thủy nông”, với “Mo cơm, quả cà và tấm lòng cộng sản !”…
Nhưng xin đừng hiểu lầm rằng Toán học chỉ có ngành Toán học kiến thiết. Và cũng đừng hiểu lầm rằng nhà Toán học Phan Đình Diệu chỉ có nghiên cứu Toán học kiến thiết !
Ngay từ cuối những năm 60 của thế kỷ trước, tại Hội trường Uỷ ban Khoa học và kỹ thuật nhà nước ở 39 Trần Hưng Đạo Hà Nội, trong một chuỗi dài các Seminar dành cho các nhà Toán học của GS Tạ Quang Bửu như về “ Các cấu trúc của Bourbaki “ hay về “ Lý thuyết các Tai biến”… , nhiều buổi GS Tạ Quang Bửu và Giảng viên Phan Đình Diệu cùng nhau thuyết trình, hoặc có buổi chỉ mình Phan Đình Diệu thuyết trình, mặc dù lúc ấy ông chưa phải là Tiến sĩ, càng chưa phải Giáo sư !
Tôi cũng như nhiều bạn bè cùng lứa, may mắn được là bạn học cùng ông từ hồi phổ thông cho đến đại học, (cũng có cái không may là phải chịu đựng cái chủ nghĩa thành phần trong một giai đoạn, nhưng rồi cũng đã qua đi). Điều mấy mắn lớn lao của thế hệ chúng tôi là ở chỗ chúng tôi đã được hưởng thụ sự giáo dục từ những người thầy trí tuệ và tâm huyết mà những thế hệ sau không dễ gì có được ! 

Có thể nói hệ thống tư duy minh bạch, trung thực,  cùng với tài năng vượt trội của Phan Đình Diệu đã tạo nên uy tín lớn lao của ông trong lòng  bạn bè, đồng nghiệp, học trò, và nhiều người khác.

Còn nhớ ngày họp Đại hội thành lập Hội Tin học Việt Nam, tổ chức tại Cung Văn hoá Hữu nghị Việt Xô (nay là Cung Hữu nghị), ông được tín nhiệm gần như tuyệt đối, mặc dù mọi người cũng đều biết rằng bên cạnh ông , còn có nhiều nhà Tin học khác, cũng tài năng và tâm huyết !  Sau khi Ban Kiểm phiếu công bố kết quả bầu Chủ tịch Hội, cả Hội trường vỗ tay vang dội, rồi ùa lên sân khấu ôm nhau chúc mừng, và từ một góc nào đó trên sân khấu đã vang lên những bài hát nhạc chế Việt và Nga quen thuộc từ thời sinh viên….

Sau này, do công việc, tôi lại có dịp làm việc với ông trong vài trò là Uỷ viên thường trực của Chương trình Tin học hoá Quản lý ở Bộ Công nghiệp nhẹ (trước đây), trong khi ông là Phó Chủ nhiệm thường trực Ban chỉ đạo Chương trình Tin học hoá Quản lý của Nhà nước. Một lần tôi đã tâm sự với ông : “Tầm nhận thức của Diệu thì rất cao, mà thực tế thì đang rất thấp. Phải có nhiều người như bọn tớ nữa để có thể chuyển tải được điều Diệu mong muốn đến với thực tế !” Ông bảo ông không đồng ý như thế, nhưng ông đã nhiều lần cùng chúng tôi bàn cách đưa Tin học vào cuộc sống như thế nào. Những khái niệm như Tin học văn phòng, Tin học Tài chính Ngân hàng, Tin học Hàng không…ở nước ta dường như đã ra đời từ thực tiễn hoạt động ấy. Nhiều người có tài năng và tâm huyết đã đóng góp vào những thành công ấy, nhưng dấu ấn của Phan Đình Diệu là không hề nhỏ !
Bây giờ thì những dấu ấn ấy đã trở thành dĩ vãng xa vời ! 

Ước ao rằng Diệu sẽ gặp Cương và các bậc thầy đáng kính của chúng ta: các GS Tạ Quang Bửu, Lê Văn Thiêm, Trần Đại Nghĩa …tại một không gian “ánh xạ” nào đó !

Vô cùng thương nhớ Diệu !

Bạn cũ : Trần Văn Kinh